\SourceLeaf{592}schung das specifische Gewicht 1 zu verschaffen. Übrigens hat diese Unklarheit hier weiter keinen Einfluss. Die Rechnung ist aber nicht sehr genau geführt, da aus den angegebenen Datis das specifische Gewicht des Sauerstoffgases nicht $1.1148$, sondern $1.11447$ folgt. Dieser Fehler. ist jedoch ganz unerheblich.

Die dritte Columne gibt die in 100 Theilen trockner Luft \textit{dem Gewichte nach} enthaltenen Theile der einzelnen Gasarten, nemlich
\begin{center}\begin{tabular}{lcrl}
Stickluft & .\quad .\quad . & $76.49$ & Theile\\
Sauerstoffluft & . & $23.41$ & —\\
Kohlensaure Luft & & $0.10$ & —
\end{tabular}\end{center}
Diese Zahlen sind offenbar nur die Producte der Zahlen der ersten Columne in die dazu gehörigen der zweiten.

Die vierte Columne gibt die Höhe an, auf welcher (unter Voraussetzung der Richtigkeit der \textsc{Dalton}schen Hypothese) jede einzelne Atmosphäre das Barometer halten würde, oder den Antheil an dem Totaldruck, letztern für trockne Luft zu $27.76$ Zoll Quecksilberhöhe angenommen. Herrn \textsc{Benzenbergs} Zahlen sind
\begin{center}\begin{tabular}{llcrl}
für die & Stickluft-Atmosphäre & .\quad .\quad . & $21.2326$ & Zoll\\
—\quad — & Sauerstoffluft-Atmosphäre & . & $6.4986$ & —\\
—\quad — & kohlensaure Luft-Atmosphäre & & $0.0278$ & —
\end{tabular}\end{center}
Wir werden auf die Berechnung dieser Columne sogleich zurückkommen.

Die fünfte Columne enthält die specifischen Gewichte der Luftarten mit Quecksilber verglichen, beim Gefrierpunkte und 28 Zoll Quecksilberdruck. Diese Zahlen, nemlich resp.\ $\frac{1}{10830}$, $\frac{1}{9414}$, $\frac{1}{6997}$ sind nichts anderes, als die Producte der Zahlen der zweiten Columne in $\frac{1}{10495}$.

Die sechste Columne hat nur die Überschrift: \textit{Beständige Zahl.} Man sieht aber, dass die sogenannte Subtangente gemeint ist, oder die Höhe, welche eine fingirte Atmosphäre von gleichförmiger und zwar so grosser Dichtigkeit, wie die wirkliche unten hat, haben müsste, um eben so stark zu drücken, wie diese. Für gemeine trockne Luft ist also diese Zahl das Product aus 10495 in 28 Zoll, oder $24488\frac{1}{3}$ Fuss; für die drei einzelnen Atmosphären, in \textsc{Daltons} Vorstellungsweise, werden die Zahlen eben so die Producte aus 28 Zoll in die Nenner der Brüche der fünften Columne sein, oder einfacher, man findet sie, wenn man $24488.33$ Fuss mit den Zahlen der zweiten Columne dividirt. Wir schreiben diese Zahlen
\SourceLeaf{593}sowohl wie sie Herr \textsc{Benzenberg} angibt, als wie sie aus einer schärfern Rechnung folgen, hier her
\begin{center}\begin{tabular}{lcrlrl}
& & \multicolumn{2}{c}{nach Hr.\ \textsc{Benz}.} & \multicolumn{2}{c}{nach schärf.\ R.}\\
Stickstoffluft & .\quad . & $25270$ & Fuss & $25269.15$ & Fuss\\
Sauerstoffluft & . & $21966$ & — & $21973.01$ & —\\
Kohlensaure Luft & & $16326$ & — & $16325.56$ & —
\end{tabular}\end{center}

Alles bisher gegen Herrn \textsc{Benzenberg} erinnerte ist durchaus unerheblich: wir kommen aber jetzt zu dem wesentlichen Punkte, der Berechnungsart der vierten Columne. Herr \textsc{Benzenberg} erklärt sich gar nicht darüber, \textit{wie} er diese Berechnung gemacht habe: er sagt blos, dass es Beispiele aus der Gesellschaftsrechnung seien. Man erkennt aber leicht, dass er die Zahlen der vierten Columne denen der dritten schlechthin proportional gesetzt, oder jene aus der Multiplication von $27.76$ Zoll mit
\begin{center}\begin{tabular}{rl}
$0.7649$ & für Stickstoffluft\\
$0.2341$ & für Sauerstoffluft\\
$0.0010$ & für kohlensaure Luft
\end{tabular}\end{center}
abgeleitet hat.

Und dies ist unrichtig.

Denn der ganze Druck der Stickstoffluft-Atmosphäre wird sich, in \textsc{Daltons} Hypothese, zu dem ganzen Druck der Sauerstoffluft-Atmosphäre, nicht wie die Gewichtsantheile, welche diese Gasarten an dem untersten Cubikfuss gemischter Luft haben, verhalten, sondern im \textit{zusammengesetzten} Verhältniss dieser Gewichtsantheile einerseits, und der den beiden Gasarten zukommenden Subtangenten andererseits, stehen, also den \textit{Producten} aus den Zahlen der dritten und sechsten Columne proportional sein müssen, oder was dasselbe ist, den Quotienten, wenn die Zahlen der dritten Columne mit denen der zweiten dividirt werden, also schlechthin den Zahlen der ersten Columne.

Bei einiger Überlegung ist dies auch von selbst klar, denn die Bedeutung der Zahlen der ersten Columne kann auch so ausgesprochen werden: die in einem Volumen von 100 Theilen gemeiner trockner Luft am Boden der Atmosphäre enthaltene Stickluft würde, für sich allein genommen, unter demselben Druck, unter welchem jene steht, nur den Raum von $76.93$ Theilen einnehmen, und eben
\SourceLeaf{594}so die andern Gasarten: indem also jede dieser drei Gasarten jetzt in den Raum von 100 Theilen verbreitet ist und von den übrigen unabhängig gedacht wird, verhalten sich die Quecksilberdrucke, denen sie einzeln das Gleichgewicht halten, wie
\begin{center}\begin{tabular}{rl}
$78.93$ & für Stickstoffluft\\
$21.00$ & für Sauerstoffluft\\
$0.07$ & für kohlensaure Luft.
\end{tabular}\end{center}

Man sieht also, dass Herr \textsc{Benzenberg} in seiner vierten Columne den Totaldruck von $27.76$ Zoll unrichtig vertheilt hat. Die richtigen Zahlen sind
\begin{center}\begin{tabular}{rlll}
$21.9110$ & Zoll & für & Stickstoffluft\\
$5.8296$ & — & — & Sauerstoffluft\\
$0.0194$ & — & — & kohlensaure Luft.
\end{tabular}\end{center}
Die Zahlen der vierten und sechsten Columne sind aber die Elemente, nach denen der Druck der Atmosphäre in jeder Höhe, in der \textsc{Dalton}schen Hypothese berechnet werden muss. Natürlich ergeben sich daher mit den verbesserten Werthen andere Resultate.

Wir haben uns die Mühe gegeben, diese Rechnung für einige Höhen zu führen. Folgende Tafel enthält die Resultate:
\begin{center}\begin{tabular}{r|r|r|r}
& \multicolumn{3}{c}{Barometer-Höhe in Zollen}\\
\multicolumn{1}{c|}{Höhe} & \multicolumn{1}{c|}{n.\ gewöhnl.} & \multicolumn{2}{c}{in \textsc{Daltons} Hypothese}\\
\multicolumn{1}{c|}{Fuss} & \multicolumn{1}{c|}{Theorie} & \multicolumn{1}{c|}{unsre R.} & \multicolumn{1}{c}{H.\ \textsc{Benz}.\ R.}\\\hline
$5000$ & $22.6332$ & $22.6350$ & $22.6179$\\
$10000$ & $18.4532$ & $18.4589$ & $18.4314$\\
$15000$ & $15.0452$ & $15.0555$ & $15.0221$\\
$20000$ & $12.2666$ & $12.2814$ & $12.2458$
\end{tabular}\end{center}

Die Zahlen für die Barometerhöhe nach der gewöhnlichen Theorie haben wir hier nach unserer eigenen Berechnung angesetzt; die von Hrn.\ \textsc{Benzenberg} für dieselben Höhen auch nach der gewöhnlichen Theorie berechneten (S.\ 12) weichen davon zum Theil etwas ab. Die Richtigkeit unserer Rechnung können wir verbürgen.

Die Barometerhöhe in der \textsc{Dalton}schen Hypothese weicht also von der Barometerhöhe nach der gewöhnlichen Theorie um folgende Unterschiede ab:
\SourceLeaf{595}\begin{center}\begin{tabular}{rl|rl|rl}
& & \multicolumn{4}{c}{Unterschied der Barometerhöhe}\\
\multicolumn{2}{c|}{Höhe} & \multicolumn{2}{c|}{nach unserer R.} & \multicolumn{2}{c}{nach H.\ \textsc{Benz}.}\\\hline
$5000$ & Fuss & $+0.0018$ & Zoll & $-0.0153$ & Zoll\\
$10000$ & — & $+0.0057$ & — & $-0.0218$ & —\\
$15000$ & — & $+0.0103$ & — & $-0.0231$ & —\\
$20000$ & — & $+0.0148$ & — & $-0.0208$ & —
\end{tabular}\end{center}

Einer bestimmten Höhe entspricht daher, in \textsc{Daltons} Hypothese, nicht, wie Herr \textsc{Benzenberg} meint, ein kleinerer, sondern ein grösserer Barometerstand, als in der gewöhnlichen Theorie, und eben so wird folglich aus einem bestimmten Barometerstande in jener Hypothese nicht eine kleinere, sondern eine grössere Höhe berechnet werden. Für den Monte Gregorio ist dieser Unterschied nicht $-16$ Fuss, sondern $+2$ Fuss. Bei kleineren Höhen wird der Unterschied sehr nahe dem Quadrat der Höhe proportional; dies liess sich auch leicht durch eine nahe liegende Betrachtung \textit{a priori} voraussehen, welche wir jedoch Kürze halber hier nicht weiter entwickeln. Herrn \textsc{Benzenbergs} Unterschiede hingegen sind für kleine Höhen diesen nahe proportional, was allein schon zureichte, die Unrichtigkeit derselben zu erkennen. Übrigens ist obige numerische Rechnung nur für trockene Luft geführt: die Berücksichtigung der Wasserdämpfe würde die relativen Unterschiede nur sehr wenig ändern.

Wir fügen dieser Darstellung noch einige Bemerkungen bei.

I.\quad Unser Resultat, dass der Unterschied der Barometerhöhe in \textsc{Daltons} Hypothese von der auf gewöhnliche Weise berechneten \textit{positiv}, und für mässige Höhen deren Quadraten nahe proportional wird, ist allgemein, und von den angenommenen Werthen der specifischen Gewichte der einzelnen Gasarten, aus denen die gemischte Luft besteht, unabhängig. Es würde also vergeblich sein, von andern Werthen dieser specifischen Gewichte ein günstigeres Resultat zu erwarten.

II.\quad Schon im Jahre 1807 hat \textsc{Tralles} eine richtige Darstellung der Berechnung der Barometerhöhen in \textsc{Daltons} Hypothese geliefert, welche man nur oberflächlich anzusehen braucht, um zu erkennen, dass sein Resultat mit dem unsrigen im Wesentlichen ganz übereinstimmt. Man muss sich daher wundern dass der Verfasser des einen oben erwähnten Artikels des physikalischen Wörterbuchs behauptet, \textsc{Tralles} finde sehr nahe dieselbe Differenz, wie Herr \textsc{Benzenberg}. In der That ist sie im Zeichen und im Gesetz verschieden (für mässige
\SourceLeaf{596}Höhen dem Quadrate der Höhe, und nicht dieser selbst proportional). Vielleicht hat ein, doch leicht als solcher zu erkennender, Druckfehler an diesem Versehen Schuld, da in dem Aufsatz von \textsc{Tralles} (\textsc{Gilberts} Annalen B.\ 27.\ S.\ 445) einmal $\frac{B}{b+\beta}$ anstatt $\frac{b+\beta}{B}$ gesetzt ist: bei der Anwendung auf ein bestimmtes Beispiel (ebendas.\ unterste Zeile) steht aber doch dieser Bruch wieder richtig, so wie auf der folgenden Seite der Unterschied mit seinem richtigen Zeichen angesetzt ist. Noch mehr aber muss man sich wundern, dass Herr \textsc{Benzenberg} es unterliess, das Resultat seines als gründlicher Mathematiker bekannten Vorgängers mit dem seinigen, dem es ganz entgegengesetzt ist, zu vergleichen. Man würde in der That vermuthen, dass Herr \textsc{Benzenberg} diesen Aufsatz gar nicht gekannt habe, wenn er nicht desselben ausdrücklich erwähnte, obwohl nur mit der Abfertigung S.\ 15, `Herr \textsc{Tralles} hat Buchstabenrechnung angewendet. Dieses ist unnöthig. Wenn man die Vorstellung von vier Barometern hat, so kann man es mit der Regel von dreien ausführen, und man gebraucht gar keine Gelehrsamkeit'. Dieser Grundsatz, zu welchem Herr \textsc{Benzenberg} sich bei vielen Gelegenheiten — wir wollen hier nicht untersuchen, ob allemal bei den rechten — laut bekannt hat, mag übrigens für den vorliegenden Fall eingeräumt werden, und unsere Darstellung, wenn es uns gelungen ist, ihr die erforderliche Klarheit zu geben, dann selbst als Bestätigung dienen.

III.\quad Wenn es nun eine vergebliche Mühe ist, den Unterschied der barometrischen und der trigonometrischen Messung des Monte Gregorio durch \textsc{Daltons} Hypothese heben zu wollen (in welcher er sogar noch um zwei Fuss vergrössert wird), so steht es als eine entschiedene Thatsache fest, dass eine von beiden, oder beide, nicht diejenige Genauigkeit haben, welche Herr \textsc{Benzenberg} ihnen beilegen zu können glaubte. Nach unserer Meinung mögen alle drei hier in Frage kommenden Fehlerquellen ihren Antheil daran haben. Erstlich das Schwanken der gemessenen Barometerhöhen selbst. Zweitens die in der Berechnung gebrauchte Constante, welche Herr \textsc{Benzenberg} auf \textsc{Biots} Abwägung der atmosphärischen Luft gegründet hatte, und die wohl viel sicherer aus einer zweckmässigen Benutzung zahlreicher zugleich barometrisch und trigonometrisch gemessener Berghöhen bestimmt werden kann. Aber drittens mag auch die trigonometrische Messung des Monte Gregorio selbst ihren Theil zu dem Unterschiede beigetragen haben, da wir dem Urtheil des Herrn \textsc{Benzenberg} über ihre unübertreffliche Genauigkeit nicht ganz beipflichten können. Die Bestimmung der Ent\SourceLeaf{597}fernung der Spitze des Berges von dem Standpunkte am Fusse desselben gründete sich nur auf eine kleine Basis. Ihre Länge (670 Meter) scheint zwar mit vieler Sorgfalt gemessen zu sein (man brachte vier Tage damit zu); allein der ihr gegenüber stehende Winkel (nur $6^0\;14'$) wurde nicht gemessen, sondern nur aus den beiden andern geschlossen. Ein solches Verfahren erfordert immer selbst bei dem Gebrauch vortrefflicher Werkzeuge grosse Behutsamkeit: allein \textsc{Daubuissons} Winkelmessungs-Instrument, ein \textsc{Lenoir}scher Repetitionskreis von acht Zoll Durchmesser, scheint nur ein sehr mittelmässiges gewesen zu sein, da wir sehen, dass von den zehn Repetitionen, aus welchen die Winkelmessungen an jedem der beiden Standpunkte bestehen, die einzelnen Paare einmal Unterschiede für den einfachen Winkel geben, die über eine Minute gehen. Auch sagt \textsc{Daubuisson} nichts über die Beschaffenheit der von ihm zu Zielpunkten gebrauchten Signale, und es lässt sich daher nicht beurtheilen, mit welcher Schärfe sich dieselben einschneiden liessen, und ob nicht dabei eine nachtheilige Beleuchtungsphase Statt finden konnte. Ein Fehler von einer halben Minute in dem geschlossenen dritten Winkel würde die gemessene Höhe schon um sieben Fuss ändern. — Übrigens können wir dies auf sich beruhen lassen, da es für die gegenwärtige Frage ganz gleichgültig ist, wie man den Unterschied erklären will: genug, dass \textsc{Daltons} Hypothese gar nichts dazu beitragen kann.

\Needspace{22\baselineskip}\vspace{6mm}
\pdfbookmark[1]{Practische Anweisung zur Verfertigung künstlicher Magnete.}{fischer-magnete}
\begin{center}\small Göttingische gelehrte Anzeigen.\quad 1832.\quad September 10.\end{center}

\textit{Practische Anweisung zur vortheilhaften Verfertigung und Zusammenfügung künstlicher Magnete, besonders der Hufeisen, geraden Stäbe, Compass- und anderer Nadeln, so wie die neueste Entdeckung, denselben die höchste Anziehungskraft zu ertheilen, für Naturforscher, Ärzte, Seefahrer, Techniker und alle andere Arten von Metallarbeitern, von \textup{\textsc{Friedrich Fischer}}, Lehrer und practischem Techniker. Heilbronn.\ \textup{1833.} In der Classschen Buchhandlung.} (58 S.\ in 8.\ Mit zwei lithographirten Tafeln.)

Unter den mannigfaltigen Phänomenen, welche die magnetische Kraft darbietet, zieht das Tragen bedeutender Lasten durch künstliche Magnete, deren Gewicht nur einen sehr kleinen Theil derselben beträgt, vorzüglich die Bewun\SourceLeaf{598}derung der Liebhaber auf sich, während es in wissenschaftlicher Rücksicht nur ein untergeordnetes Interesse hat, und in dem reichhaltigsten Werke der neuesten Zeit über die Physik kaum mit einigen Worten erwähnt wird. Die Anordnung der für jenen Zweck am meisten geeigneten künstlichen Magnete in Hufeisenform, wird daher weniger als Sache des Physikers betrachtet, sondern ist mehr in den Händen von Personen, die einen Erwerb daraus machen, und zuweilen angeblich neue und eigenthümliche Methoden unter dem Siegel des Geheimnisses zum Verkauf ausbieten. Durch einen neuern Fall dieser Art ist der Verfasser der vorliegenden kleinen Schrift zu eigenen Versuchen veranlasst, und es gereicht ihm zur Ehre, dass er die Resultate derselben ohne Rückhalt und Geheimnisskrämerei veröffentlicht. Bei weitem der grösste Theil der Schrift ist den künstlichen Magneten in Hufeisenform gewidmet, ihre vortheilhaftesten Verhältnisse, die Auswahl und Behandlung des Stahls und die anzuwendenden Streichmethoden werden auf eine fassliche Art beschrieben, und die Liebhaber können versichert sein, dass sie durch Befolgung der gegebenen Vorschriften sich allezeit solche Magnete von sehr grossem Tragvermögen verschaffen.

Referent würde sich auf diese Empfehlung der vorliegenden Schrift beschränken, wenn nicht eben die von dem Verfasser gebrauchten Streichmethoden (die vermuthlich unter der auf dem Titel erwähnten neuesten Entdeckung verstanden sein sollen, obwohl er selbst einräumt, dass solche auch sonst bekannt sein mögen) dem Referenten zu einigen eigenen Bemerkungen Anlass gäben, die als eine nicht unwichtige Ergänzung von \textsc{Coulombs} Erfahrungen über die allen Physikern wohlbekannten Methoden von \textsc{Knight}, \textsc{Duhamel}, \textsc{Michell}, \textsc{Canton} und \textsc{Aepinus} betrachtet werden können.

Der Verfasser bedient sich zur Erregung des Magnetismus in einem anzufertigenden Hufeisenmagnet eines schon vorhandenen Magnets von derselben Form, und sein Verfahren besteht aus zwei nach einander anzuwendenden Operationen, wo immer beide Pole zugleich streichen, aber in der ersten der eine Pol des Streichmagnets dem andern auf seinem Wege folgt, in der zweiten hingegen der eine Pol auf dem einen Arm, der andere auf den andern von der Krümmung nach dem Ende zu geführt wird. Es ist unnöthig, die dabei erforderliche Ordnung der Pole hier besonders zu bemerken. Vor der zweiten Operation räth der Verfasser noch an, den zu bestreichenden Magnet zu erwärmen, und die Arbeit bis zu erfolgter Abkühlung fortzusetzen.

\SourceLeaf{599}Man sieht nun leicht, dass die erste Operation mit dem vom \textsc{Michell} erfundenen Doppelstrich ganz einerlei ist. Die zweite Operation kommt hingegen im Wesentlichen mit \textsc{Duhamels} Verfahren überein, nur dass die von \textsc{Duhamel} zum Streichen angewandten getrennten geraden Stäbe (oder Büschel von Stäben) einige Vortheile für kräftigere Erregung gewähren, deren man bei Anwendung Eines Hufeisen-Magnets entbehrt (besonders insofern man nicht von der Mitte der Krümmung ausgehen kann). Da nun bekanntlich \textsc{Cantons} Methode lediglich in einer Verbindung der Methoden von \textsc{Michell} und \textsc{Duhamel} besteht, so ist das Verfahren des Verfassers im Wesentlichen nur das \textsc{Canton}sche mit den Modificationen, die die Anwendung eines hufeisenförmigen Streichmagnets von selbst mit sich bringt, und enthält daher nichts eigentlich Neues, als die vorgängige Erwärmung, deren Wirksamkeit jedoch wohl erst noch weiterer Bewährung bedürfen wird: Referent hat in einigen von ihm angestellten Versuchen gar keine besondere Wirkung davon gefunden.

Was nun aber hier besonders bemerkt werden muss, ist der Umstand, dass die Physiker, nach \textsc{Coulombs} Vorgange, die Methode von \textsc{Canton} gar nicht als eine Verbesserung gelten lassen, weil, nach dem Urtheil jenes berühmten Physikers, immer nur die zuletzt angewandte Methode die Intensität des erregten Magnetismus bestimme, und daher das Vorangehen von \textsc{Michells} Streichart etwas ganz Überflüssiges sei. Von der andern Seite sieht man aus den Äusserungen unsers Verfassers, dass er die Vereinigung seiner beiden Operationen als wesentlich betrachtet, und Referent erkennt gern an, dass er selbst durch diese Äusserungen, die das Gepräge anspruchsloser Wahrheitsliebe tragen, zuerst veranlasst wurde, die Allgemeingültigkeit des Princips, welches \textsc{Coulombs} Urtheil zum Grunde liegt, in Zweifel zu ziehen: eine zahlreiche Menge von Versuchen, bei denen eigenthümliche, die grösste Schärfe gewährende, an einem andern Orte zu beschreibende Prüfungsmittel angewendet wurden, haben diesen Zweifel vollkommen gerechtfertigt.

Bekanntlich hat diejenige Verbesserung von \textsc{Michells} Streichmethode, welche wir \textsc{Aepinus} verdanken, die ausgezeichnetste Wirksamkeit, so dass bei etwas stärkern Stählen jede andere, und auch die \textsc{Duhamel}sche, bedeutend gegen sie zurücksteht. \textsc{Coulombs} Versuche haben dies ausser allen Zweifel gesetzt, und die Physiker gebrauchen daher zur kräftigsten Magnetisirung solcher Stähle ausschliesslich die Methode von \textsc{Aepinus}. Merkwürdig, und nach den bisher ange\SourceLeaf{600}nommenen Voraussetzungen unerwartet ist daher das Resultat, welches aus den erwähnten Versuchen des Referenten übereinstimmend hervorgegangen ist, dass die nach \textsc{Aepinus} Methode so stark wie möglich magnetisirten Stähle allemal noch einen bedeutenden Zuwachs von Kraft erhalten, wenn sie \textit{nachher} noch wiederholt nach \textsc{Duhamels} Verfahren gestrichen werden, wenn gleich letzteres für sich allein nur eine bedeutend schwächere Kraft entwickeln kann, als \textsc{Aepinus} Methode. Referent begnügt sich hier, diese Thatsache anzuzeigen, ohne in den Versuch einer übrigens ziemlich nahe liegenden Erklärung einzugehen. Obgleich diese Erfahrungen unmittelbar nur an der Magnetisirung gerader Stäbe gemacht sind, so ist doch nicht zu zweifeln, dass die Verbindung von \textsc{Aepinus} und \textsc{Duhamels} Methode eben so auch in hufeisenförmigen Lamellen die möglich stärkste Entwickelung des Magnetismus hervorbringen muss, nur erfordert dann die Anwendung derselben in ihrer Reinheit, wenn sie mit Bequemlichkeit ausgeführt werden soll, einige besondere Vorkehrungen. Wer diese nicht treffen mag, oder passende gerade Stäbe nicht zur Hand hat, wird, wenn auch bei etwas dickern Lamellen, nicht die höchste erreichbare, doch immer eine sehr grosse Stärke erhalten, wenn er nach des Verfassers Vorschrift einen hufeisenförmigen Streichmagnet anwendet, dessen Handhabung zugleich mit aller Bequemlichkeit geschieht.

Was der Verfasser von der Magnetisirung gerader Stäbe sagt, beschränkt sich auf die Manipulationen, die man anzuwenden hat, wenn man die Bestreichung mit einem Hufeisenmagnet ausführen will. Man erhält dadurch zwar eine grosse, aber nicht eine eben so grosse Stärke, wie durch die oben erwähnte Folge von \textsc{Aepinus} und \textsc{Duhamels} Methoden, die auch in Rücksicht auf Bequemlichkeit nichts zu wünschen übrig lassen.

Die Art, wie der Verfasser magnetisirte gerade Stäbe aufzubewahren empfiehlt, nemlich sie mit den gleichnamigen Polen auf einander zu legen, ist ganz verwerflich, wenn man wünscht, dass sie so viel wie möglich ihre Kraft behalten sollen. Am besten ist es, sie paarweise in geringer Entfernung so neben einander zu legen, dass ungleichnamige Pole zusammenkommen, und Anker aus ganz weichem Eisen von schicklicher Länge daran zu legen.

\SourceLeaf{601}\Needspace{22\baselineskip}\vspace{6mm}
\pdfbookmark[1]{Resultate aus dem Beobachtungen des magnetischen Vereins im Jahre 1836.}{resultate-1836-anzeige}
\begin{center}\small Gottingische gelehrte Anzeigen.\quad 1837.\quad Junius 29.\end{center}

\textit{Resultate aus dem Beobachtungen des magnetischen Vereins im Jahre \textup{1836.} Herausgegeben von \textup{\textsc{Carl Friedrich Gauss}} und \textup{\textsc{Wilhelm Weber}}. Göttingen \textup{1838.} Im Verlage der Dieterichschen Buchhandlung.} (124 Seiten in 8., nebst 10 Steindrucktafeln.)

Durch den Titel dieses Werks wird nur ein Theil des Inhalts bezeichnet, derjenige nemlich, welcher die nächste Veranlassung dazu gegeben hat. Von dem Vereine, welcher sich seit mehreren Jahren gebildet hat, um diejenigen Erscheinungen des tellurischen Magnetismus, die zu den interessantesten gehören, in bestimmten verabredeten Terminen gleichzeitig zu beobachten, ist schon mehrere Male in diesen Blättern die Rede gewesen (1834.\ Aug.\ 9.; 1835.\ März 7.\ [S.\ 518.\ 529 d.\ B.]), und es ist daher unnöthig, Bekanntes hier zu wiederholen. Die Theilnahme an diesem Vereine befasst schon eine grosse Anzahl von Örtern innerhalb und ausserhalb Deutschlands, und ist fortwährend im Zunehmen begriffen. Die Mittheilung der immer reichhaltiger werdenden Resultate konnte nicht länger auf einen Privatverkehr durch Briefwechsel beschränkt bleiben, sondern eine Veröffentlichung durch den Druck wurde ein Bedürfniss nicht blos für die unmittelbaren Theilnehmer, sondern auch deshalb, damit die Resultate ein Gemeingut Aller werden können, die ein Interesse an den Naturwissenschaften nehmen. Zugleich aber bietet die Herausgabe die angemessenste Gelegenheit dar, um in besondern damit zu verbindenden kleinern und grössern Aufsätzen nach und nach zur Sprache zu bringen, nicht allein was in unmittelbarem Zusammenhange mit dem nächsten Gegenstande steht, oder in mittelbarem, wie die Instrumente, ihre Berichtigung und Behandlung. sondern auch Anderes, was nur immer zu strengerer wissenschaftlicher Begründung der Lehre vom Magnetismus und Galvanismus beitragen kann.

Die vorliegende erste Lieferung enthält die graphischen Darstellungen der magnetischen Variationsbeobachtungen von sechs Terminen auf eben so vielen Steindrucktafeln, zusammen 46 Curven aus vierzehn verschiedenen Beobachtungsörtern, auch von den drei letzten Terminen die Beobachtungen selbst in Zahlen. Den grösseren Theil des Werks machen aber ausser einer historischen Einleitung
